\documentclass{article}

% Ćirilica: kodiranje fonta i jezik (srpski, ćirilica)
\usepackage[T2A]{fontenc}
\usepackage[utf8]{inputenc}
\usepackage[serbianc]{babel}

% Veličina strane i margine
\usepackage[a4paper,top=2cm,bottom=2cm,left=3cm,right=3cm,marginparwidth=1.75cm]{geometry}

% Korisni paketi
\usepackage{amsmath}
\usepackage{graphicx}
\usepackage{float}
\usepackage{booktabs}
\usepackage{tcolorbox}
\usepackage{enumitem}
\tcbuselibrary{breakable}
\usepackage[colorlinks=true, allcolors=blue]{hyperref}

% Цртежи (TikZ)
\usepackage{tikz}
\usetikzlibrary{arrows.meta, calc, decorations.markings, babel}

% TikZ стилови који се понављају на цртежима
\tikzset{
    pomocna/.style={dashed, gray!70, thin},
    mreza/.style={gray!25, very thin},
    kriva/.style={very thick, red!75!black},
}

% Оквир за решене примере
\newtcolorbox[auto counter]{primer}{breakable,
    colback=white, colframe=black!55, colbacktitle=black!55, coltitle=white,
    fonttitle=\bfseries\small, title={Пример~\thetcbcounter},
    boxrule=0.5pt, arc=1.5pt, left=6pt, right=6pt, top=3pt, bottom=3pt,
    before skip=8pt, after skip=10pt}

% Оквир за кључне формуле
\newtcolorbox{kljucno}{colback=blue!4, colframe=blue!45!black,
    boxrule=0.5pt, arc=1.5pt, left=8pt, right=8pt, top=4pt, bottom=4pt}

\title{ОЕТ 2 -- Комплексна импеданса и адмитанса}
\author{Саобраћајно техничка школа, Београд, Земун}
\date{}

\begin{document}
\maketitle

%==================================================================
\section{Бројеви $1$, $j$, $-1$ и $-j$ на јединичном кругу}
%==================================================================

По Ојлеровој формули, тачка на јединичном кругу под углом $\theta$ је
$e^{j\theta} = \cos\theta + j\sin\theta$. Ако уврстимо четири карактеристична
угла, добијамо четири тачке у којима круг сече осе:

\begin{center}
\renewcommand{\arraystretch}{1.3}
\begin{tabular}{lcccc}
    \toprule
    $\theta$                    & $0^\circ$      & $90^\circ$     & $180^\circ$     & $270^\circ$ \\
    $\cos\theta + j\sin\theta$  & $1 + j\cdot 0$ & $0 + j\cdot 1$ & $-1 + j\cdot 0$ & $0 + j\cdot(-1)$ \\
    резултат                    & $1$            & $j$            & $-1$            & $-j$ \\
    \bottomrule
\end{tabular}
\end{center}

Сва четири броја имају модуо 1, па при множењу \textbf{не мењају дужину}
фазора, већ само \textbf{додају свој угао}. Зато су они, у ствари, упутства
за окретање:

\begin{kljucno}
\begin{tabular}{ll}
    $\cdot\,1$    & ништа се не мења \\
    $\cdot\,j$    & окрет за $+90^\circ$ (супротно од казаљке): сигнал \textbf{претходи} за $90^\circ$ \\
    $\cdot\,(-1)$ & окрет за $180^\circ$: сигнал се \textbf{обрне} (промени знак) \\
    $\cdot\,(-j)$ & окрет за $-90^\circ$ (у смеру казаљке): сигнал \textbf{касни} за $90^\circ$ \\
\end{tabular}
\end{kljucno}

Ако напону $u = U_m\sin\omega t$ одговара фазор $\underline{U}$, фазору
$j\underline{U}$ одговара $U_m\sin(\omega t + 90^\circ)$, а фазору $-j\underline{U}$
одговара $U_m\sin(\omega t - 90^\circ)$. Види се и зашто је $j^2 = -1$: два окрета
од $90^\circ$ дају окрет од $180^\circ$.

\begin{figure}[H]
    \centering
    \begin{tikzpicture}[>=Stealth, font=\small]
        % --- фазори ---
        \draw[->, gray] (-1.5,0) -- (1.6,0) node[right] {Re};
        \draw[->, gray] (0,-1.5) -- (0,1.6) node[above] {Im};
        \draw[->, very thick, blue!70!black]  (0,0) -- (1.1,0)  node[below right] {$\underline{U}$};
        \draw[->, very thick, red!75!black]   (0,0) -- (0,1.1)  node[right] {$j\underline{U}$};
        \draw[->, very thick, gray!60!black]  (0,0) -- (-1.1,0) node[below left] {$-\underline{U}$};
        \draw[->, very thick, green!50!black] (0,0) -- (0,-1.1) node[right] {$-j\underline{U}$};
        % --- временски дијаграм ---
        \begin{scope}[xshift=2.9cm]
            \def\s{0.0125}
            \draw[->] (0,0) -- (4.9,0) node[right] {$\omega t$};
            \draw[->] (0,-1.2) -- (0,1.3);
            \foreach \d/\l in {90/{90^\circ},180/{180^\circ},270/{270^\circ},360/{360^\circ}}
                \node[below, font=\scriptsize] at ({\d*\s},-1.0) {$\l$};
            \draw[very thick, blue!70!black, domain=0:360, samples=80] plot ({\x*\s},{0.9*sin(\x)});
            \draw[very thick, red!75!black, domain=0:360, samples=80] plot ({\x*\s},{0.9*sin(\x+90)});
            \draw[very thick, green!50!black, domain=0:360, samples=80] plot ({\x*\s},{0.9*sin(\x-90)});
            \node[blue!70!black, right, font=\scriptsize] at (5.6,0.7) {$u$};
            \node[red!75!black, right, font=\scriptsize] at (5.6,0.3) {$j\cdot u$: претходи};
            \node[green!50!black, right, font=\scriptsize] at (5.6,-0.1) {$-j\cdot u$: касни};
        \end{scope}
    \end{tikzpicture}
    \caption{Множење са $j$, $-1$ и $-j$ окреће фазор (лево) и помера синусоиду (десно)}
    \label{fig:j-okret}
\end{figure}

%==================================================================
\section{Импеданса}
%==================================================================

У колу наизменичне струје отпорник, калем и кондензатор се супротстављају
струји, али калем и кондензатор уз то и \textbf{померају фазу} струје за
$90^\circ$. \textbf{Импеданса} $\underline{Z}$ је комплексна отпорност која у
једном броју носи обе информације: колико елемент смањује струју и за колико
је помера.

\begin{kljucno}
\[
\underline{Z} = \frac{\underline{U}}{\underline{I}} = R + jX = |\underline{Z}|\,e^{j\varphi}
\qquad [\Omega],
\qquad
|\underline{Z}| = \sqrt{R^2 + X^2}, \quad \varphi = \operatorname{arctg}\frac{X}{R}
\]
\end{kljucno}

Реални део $R$ је \textbf{резистанса}: она троши енергију и претвара је у
топлоту. Имагинарни део $X$ је \textbf{реактанса}: она енергију не троши, већ је
привремено чува у пољу и враћа извору. За основне елементе ($X_L = \omega L$, $X_C = \frac{1}{\omega C}$, и
$\frac{1}{j} = \frac{j}{j^2} = -j$):
\[
\underline{Z}_R = R, \qquad
\underline{Z}_L = j\omega L = jX_L, \qquad
\underline{Z}_C = \frac{1}{j\omega C} = -jX_C
\]

\textbf{Зашто нам треба?} Са импедансом за наизменичну струју важе
\textbf{исти закони} као за једносмерну: Омов закон
$\underline{U} = \underline{Z}\,\underline{I}$, Кирхофови закони, редна веза
$\underline{Z}_e = \underline{Z}_1 + \underline{Z}_2$ и паралелна веза. Једина
разлика је што рачунамо са комплексним бројевима, уместо да решавамо сложене
једначине са синусима.

%==================================================================
\section{Адмитанса}
%==================================================================

\textbf{Адмитанса} је реципрочна вредност импедансе, као што је проводност
реципрочна вредност отпорности: $\underline{Y} = \dfrac{1}{\underline{Z}} = G + jB$~[S].
$G$ је \textbf{кондуктанса}, а $B$ \textbf{сусцептанса}. Адмитансу рачунамо
тако што бројилац и именилац помножимо конјугованим имениоцем:
\[
\underline{Y} = \frac{1}{R + jX} = \frac{R - jX}{(R + jX)(R - jX)} = \frac{R - jX}{R^2 + X^2}
\qquad\text{и}\qquad
\underline{Y}_R = \frac{1}{R}, \quad
\underline{Y}_L = \frac{-j}{\omega L}, \quad
\underline{Y}_C = j\omega C
\]
Адмитанса је згодна код \textbf{паралелне везе}, јер се адмитансе грана
сабирају: $\underline{Y}_e = \underline{Y}_1 + \underline{Y}_2$.

\textbf{Адмитанса} $\underline{Y}$ је реципрочна вредност импедансе, као што
је проводност реципрочна вредност отпорности. Показује колико лако елемент
пропушта наизменичну струју.

\begin{kljucno}
\[
\underline{Y} = \frac{1}{\underline{Z}} = \frac{\underline{I}}{\underline{U}} = G + jB = |\underline{Y}|\,e^{-j\varphi}
\qquad [\text{S}],
\qquad
|\underline{Y}| = \sqrt{G^2 + B^2} = \frac{1}{|\underline{Z}|}
\]
\end{kljucno}

\begin{itemize}[noitemsep]
    \item $\underline{Y}$ -- \textbf{адмитанса} (комплексна проводност),
    \item $G$ -- \textbf{кондуктанса} (реални део, активна проводност), одговара резистанси,
    \item $B$ -- \textbf{сусцептанса} (имагинарни део, реактивна проводност), одговара реактанси.
\end{itemize}
Јединица за све три величине је \textbf{сименс} (S), $1~\text{S} = 1/\Omega$.

Адмитансу рачунамо тако што бројилац и именилац помножимо конјугованим
имениоцем:
\[
\underline{Y} = \frac{1}{R + jX} = \frac{R - jX}{(R + jX)(R - jX)} = \frac{R - jX}{R^2 + X^2}
\quad\Longrightarrow\quad
G = \frac{R}{R^2 + X^2}, \qquad B = \frac{-X}{R^2 + X^2}
\]

\textbf{Сусцептансе калема и кондензатора:}
\[
B_L = \frac{1}{\omega L} = \frac{1}{X_L} \quad\text{(индуктивна)}, \qquad
B_C = \omega C = \frac{1}{X_C} \quad\text{(капацитивна)}
\]
па су адмитансе основних елемената:
\[
\underline{Y}_R = G = \frac{1}{R}, \qquad
\underline{Y}_L = \frac{1}{j\omega L} = -jB_L, \qquad
\underline{Y}_C = j\omega C = jB_C
\]

Адмитанса је згодна код \textbf{паралелне везе}, јер се адмитансе грана
сабирају: $\underline{Y}_e = \underline{Y}_1 + \underline{Y}_2$, а кондуктансе и
сусцептансе посебно: $G_e = G_1 + G_2$, $B_e = B_1 + B_2$.

\begin{center}
\renewcommand{\arraystretch}{1.3}
\begin{tabular}{lll}
    \toprule
    \textbf{Импеданса} $\underline{Z}$ [$\Omega$] & \textbf{Адмитанса} $\underline{Y}$ [S] & \\
    \midrule
    резистанса $R$  & кондуктанса $G$  & реални део \\
    реактанса $X$   & сусцептанса $B$  & имагинарни део \\
    $X_L = \omega L$ & $B_L = 1/(\omega L)$ & калем \\
    $X_C = 1/(\omega C)$ & $B_C = \omega C$ & кондензатор \\
    \bottomrule
\end{tabular}
\end{center}

\end{document}